Le colorant (\textit{dye} en anglais) est représenté par un champ
scalaire de concentration $c(\mathbf{x}, t)$, transporté par
advection dans le champ de vitesse du fluide.
Il constitue un traceur passif : sa présence n'influence pas
la dynamique de l'écoulement, mais permet d'en visualiser le transport.

En l'absence de diffusion, l'évolution de la concentration satisfait
l'équation suivante sous
\conceptlink{annexes:forme_conservative}{forme conservative} :
\begin{equation}
    \label{eq:dye_transport}
    \boxed{
        \frac{\partial c}{\partial t}
        + \nabla \cdot (c\mathbf{u}) = s.
    }
\end{equation}
Le terme source $s(\mathbf{x}, t)$ représente la quantité de
concentration ajoutée par unité de temps. Ce terme pourra être utilisé pour répondre à
l'injection de colorant par l'utilisateur.